\chapter{Project Tools \& the Example Catalog}
\label{ch:project-cli}

\newthought{Projects are the unit of sharing.} \cli{Jarvis project} scaffolds new
projects, packs existing ones into tarballs, browses the official example catalog, and
--- for restricted material --- encrypts and decrypts packs so that no user ever touches
\cli{openssl} by hand.

\section{\clih{project create}: scaffolding}
\label{sec:project-create}

\begin{terminal}
\begin{jcmd}
Jarvis project create MyScan
cd MyScan
Jarvis run bin/quickstart_bridson_operas.yaml
\end{jcmd}
\end{terminal}

\begin{jfiletree}
\dirtree{%
.1 \dtfoldopen{MyScan/}.
.2 \dtyaml{jarvis.project.yaml}\DTcomment{descriptor (root marker)}.
.2 \dtmark{.jarvis-project.json}\DTcomment{machine marker}.
.2 \dtfold{bin/}\DTcomment{task cards, incl.\ a runnable quick start}.
.2 \dtfold{data/}\DTcomment{input data, check points, CSVs}.
.2 \dtfold{deps/}\DTcomment{environment\_default.yaml (EnvReqs.V2 defaults)}.
}
\end{jfiletree}

The markers pin the project root (Section~\ref{sec:projects}); the scaffolded defaults
file wires \yamlkey{EnvReqs.Check\_default\_dependencies} so every card in \file{bin/}
shares one runtime profile (Section~\ref{sec:check-default-deps}).

\section{\clih{project pack}: tarballs}
\label{sec:project-pack}

\begin{terminal}
\begin{jcmd}
Jarvis project pack . --share    # curated: bin/ data/ outputs/ images/ + markers
Jarvis project pack . --repro    # default: all non-hidden top-level content
Jarvis project pack . --full     # everything
Jarvis project pack . --man      # write an editable pack manifest only
\end{jcmd}
\end{terminal}

\begin{keylist}
  \item[--share] The publishable subset: task cards, data, outputs, images, the
        project markers, and \file{README.md} --- no local tool installs, no scratch.
  \item[--repro] The default: every non-hidden top-level entry, i.e.\ enough to re-run
        the project elsewhere.
  \item[--full] Everything under the root.
  \item[--man] Write only a manifest \textsc{yaml} listing what \emph{would} be packed
        --- edit it, then pack from it; also your audit trail of what left the machine.
\end{keylist}

All profiles silently exclude the usual junk (\file{.git/}, \code{\_\_pycache\_\_},
caches, \code{.pyc}/\code{.tmp}/swap files, \code{.DS\_Store}).

\section{The official catalog}
\label{sec:catalog}

Curated example projects are indexed by a single \textsc{json} document in the
\code{Jarvis-Examples} repository; no separate package to install.

\begin{terminal}
\begin{jcmd}
Jarvis project list          # table: Name, Access, Key, Category, Summary
Jarvis project browse        # the same listing, under a second name
Jarvis project info Eggbox   # details, incl. entry-point card and key hint
Jarvis project fetch Eggbox  # download + unpack a public project
\end{jcmd}
\end{terminal}

The index is resolved in order: the remote \textsc{url} (or the
\code{JARVIS\_OFFICIAL\_LIBRARY\_INDEX\_URL} override) $\rightarrow$ the last
successfully pulled copy in \file{\textasciitilde/.jarvis/cache/} $\rightarrow$ a
snapshot packaged with the runtime. Offline machines therefore still \cli{list}, with a
possibly stale view.

\section{Restricted projects}
\label{sec:restricted}

Catalog entries marked \code{Access: restricted}, \code{Key: required} ship as
encrypted archives (\file{*.tar.gz.jenc}). The key travels out-of-band; the CLI does
the rest.

\paragraph{Fetching.}

\begin{terminal}
\begin{jcmd}
Jarvis project fetch SecretName --key 'YOUR_KEY'
# or set it once per shell:
export JARVIS_PROJECT_FETCH_KEY='YOUR_KEY'
Jarvis project fetch SecretName
\end{jcmd}
\end{terminal}

Fetch downloads, decrypts, and unpacks in one step. A wrong or missing key fails with
the catalog's key \emph{hint} for that entry.

\paragraph{Publishing (maintainers).}

\begin{terminal}
\begin{jcmd}
# one-shot: pack then encrypt -> MyPrivate_repro_....tar.gz.jenc
Jarvis project pack MyPrivate --repro --encrypt --key 'YOUR_KEY'

# or encrypt an existing tarball
Jarvis project encrypt MyPrivate_repro_....tar.gz --key 'YOUR_KEY'
\end{jcmd}
\end{terminal}

Upload the \file{.jenc}, register the entry in the catalog \textsc{json} with
\code{access: restricted} and \code{requires\_key: true}, and share the key privately.
Never publish the plaintext tree.

\begin{jnote}
Under the hood the format is OpenSSL-compatible \textsc{aes}-256-\textsc{cbc} with
\textsc{pbkdf2} key derivation, using the system \cli{openssl} when present (the
\code{cryptography} package otherwise). What encryption does \emph{not} hide: the
catalog metadata (names, summaries, \textsc{url}s) is public, and it protects archive
\emph{content} only --- the security model is a shared passphrase, not per-user grants.
\end{jnote}

\begin{jwarn}
Treat keys like passwords: pass them via the environment variable in shared shell
histories, never commit them, and rotate by re-encrypting and updating the catalog
entry.
\end{jwarn}
